\documentclass[../../main.tex]{subfiles}
\begin{document}
\section{Route C, normalization layer: CMH made precise, its Hodge content, and gate zero}
\label{sec:cmh-normalization}

Section~\ref{sec:moment-map-cmh} posed Question~\ref{q:cmh-normalization}: fix every object in
the schematic estimate $\norm{\Sigma^{-1/2}H\nabla g}_2^2\le4\norm{-Lg}_2^2$ and prove that the
resulting statement implies a universal Poincar\'e bound. This section answers it. The estimate
is named, its operator data are fixed, the reduction to the affine Poincar\'e inequality is
proved, and two structural consequences are recorded that were not visible while the endpoint
was a schema:

\begin{itemize}[leftmargin=1.4em]
  \item an exact weighted Hodge decomposition showing that CMH dominates the affine
        Poincar\'e constant and contains an additional solenoidal channel in dimension at least
        two; whether that channel makes CMH genuinely stronger than KLS remains open
        (Proposition~\ref{prop:cmh-hodge}, Corollary~\ref{rem:cmh-stronger-than-kls});
  \item a necessary linear-sector condition, \emph{gate zero}, which is an operator-to-trace
        upgrade of exactly the kind catalogued in Remark~\ref{rem:trace-upgrade-unification},
        and which Theorem~\ref{thm:letwin-moment-map} provably does not supply by matrix algebra
        alone (Conjecture~\ref{conj:gate-zero}, Proposition~\ref{prop:letwin-not-gate-zero}).
\end{itemize}

Throughout, $\mu$ is centered, full-dimensional and log-concave with covariance
$\Sigma=\Cov_\mu\succ0$, and $\varphi$, $\nu$, $H$ are the moment-map data of
\eqref{eq:moment-measure}--\eqref{eq:MA} in Section~\ref{sec:family-moment-map}, transported to
target coordinates as in \eqref{eq:stein-kernel-def}. Regularity is assumed only to justify
pointwise calculation; Appendix-level conventions for the weighted divergence, the closed forms,
and the affine covariance are collected in \S\ref{subsec:cmh-conventions}.

\subsection{The Stein generator and the CMH constant}
\label{subsec:cmh-definition}

Recall the affine Poincar\'e constant
\begin{equation}
  \CPaff(\mu)
  =\sup_{f\in H^1(\mu),\,f\not\equiv\mathrm{const}}
   \frac{\Var_\mu f}{\E_\mu\inner{\Sigma\nabla f}{\nabla f}},
  \label{eq:affine-poincare-constant}
\end{equation}
so that Conjecture~\ref{conj:kls} is the assertion that $\CPaff$ is bounded by a universal
constant over all dimensions and all log-concave $\mu$.

Transporting the source diffusion \eqref{eq:calL-def} to the target gives the \emph{Stein
generator}
\begin{equation}
  L_\mu g=\Div_\mu(H\nabla g)=\Tr(HD^2g)-p\cdot\nabla g,
  \qquad \Div_\mu u=\rho^{-1}\Div(\rho u),
  \label{eq:stein-generator}
\end{equation}
where $\rho$ is the density of $\mu$. With the sign convention that
$\Aop=-L_\mu\ge0$, the associated Dirichlet form is
\begin{equation}
  -\E_\mu[fL_\mu g]=\E_\mu\inner{H\nabla f}{\nabla g}.
  \label{eq:stein-dirichlet-form}
\end{equation}
That $L_\mu$ is symmetric for $\mu$ --- rather than for some auxiliary measure --- is exactly
the Stein-kernel identity $\Div_\mu H=-p$ of \eqref{eq:stein-identity}; it is what makes the
following definition canonical rather than a choice.

\begin{proposition}[Integrated Bochner identity for the Stein generator]
\label{prop:cmh-bochner}
For $g$ in the operator core,
\begin{equation}
  \E_\mu(L_\mu g)^2
  =\E_\mu\!\left[\inner{H\nabla g}{\nabla g}+\Tr(HD^2g\,HD^2g)\right].
  \label{eq:cmh-bochner}
\end{equation}
\end{proposition}

\begin{proof}
We calculate in target coordinates, use Einstein summation, and first take $g$ in a smooth
compactly supported core.  Write $\rho$ for the density of $\mu$.  The Stein identity is
\[
  \partial_i(\rho H_{ij})=-\rho p_j,
  \qquad L_\mu g=H_{k\ell}g_{k\ell}-p_k g_k.
\]
One integration by parts, followed by differentiating the displayed formula for $L_\mu g$,
gives
\begin{align*}
  \E_\mu(L_\mu g)^2
  ={}&\E_\mu[H_{ij}g_i g_j]
     -\E_\mu[H_{ij}g_j(\partial_iH_{k\ell})g_{k\ell}]\\
    &-\E_\mu[H_{ij}g_jH_{k\ell}g_{ik\ell}]
     +\E_\mu[H_{ij}g_jp_k g_{ik}].
\end{align*}
Integrating the third-derivative term by parts in $p_k$ yields
\begin{align*}
 -\E_\mu[H_{ij}g_jH_{k\ell}g_{ik\ell}]
  ={}&-\E_\mu[p_\ell H_{ij}g_jg_{i\ell}]
      +\E_\mu[H_{k\ell}(\partial_kH_{ij})g_jg_{i\ell}]\\
     &+\E_\mu[H_{k\ell}H_{ij}g_{jk}g_{i\ell}].
\end{align*}
The first term here cancels the last term in the preceding display after relabelling
$k$ and $\ell$, while the last one equals $\Tr(HD^2g\,HD^2g)$.

It remains to cancel the two terms containing $\partial H$.  If $p=\nabla\varphi(y)$, then
$H_{k\ell}(p)=\varphi_{k\ell}(y)$ and
\begin{equation}
  H_{mj}\partial_jH_{k\ell}
  =H_{mj}(H^{-1})_{rj}\varphi_{k\ell r}
  =\varphi_{mk\ell}.
  \label{eq:cmh-codazzi}
\end{equation}
Thus $H_{mj}\partial_jH_{k\ell}$ is totally symmetric in $m,k,\ell$.  Using symmetry of
$H$ and \eqref{eq:cmh-codazzi}, the residual is
\[
  \E_\mu[g_jg_{i\ell}\varphi_{\ell ij}]
  -\E_\mu[g_jg_{k\ell}\varphi_{jk\ell}]=0.
\]
This proves \eqref{eq:cmh-bochner} on the smooth core.  If $g_q$ is graph-norm Cauchy there,
apply the identity to $g_q-g_r$.  Nonnegativity of the two right-hand terms shows that
$H^{1/2}\nabla g_q$ and $H^{1/2}D^2g_qH^{1/2}$ are Cauchy in their respective $L^2$ spaces.
Their limits define the closed right-hand side, and passage to the limit proves the identity on
the operator-core closure.
\end{proof}

For $\mu$ Gaussian, $H=\Id$ and \eqref{eq:cmh-bochner} specializes to the familiar
Ornstein--Uhlenbeck identity
$\E(Lg)^2=\E[\abs{\nabla g}^2+\norm{D^2g}_{\HS}^2]$.  This is a calibration of the formula; the
proof above is analytic and uses no numerical evidence.

\begin{definition}[The canonical moment-Hessian constant]
\label{def:cmh}
The \emph{canonical moment-Hessian constant} of $\mu$ is
\begin{equation}
  \CMH(\mu)
  =\sup_{g\in\Dom(\Aop)\setminus\ker\Aop}
   \frac{\E_\mu\inner{H\nabla g}{\Sigma^{-1}H\nabla g}}{\E_\mu(L_\mu g)^2}.
  \label{eq:cmh-constant}
\end{equation}
The assertion $\CMH(\mu)\le C$ is written $\mathrm{CMH}(C)$. Here $L^2$ means $L^2(\mu)$,
$\Aop=-L_\mu$ is the nonnegative self-adjoint operator of the closed form
\eqref{eq:stein-dirichlet-form}, and every inverse is the pseudoinverse on
$(\ker\Aop)^\perp$. For a measure supported on a proper affine subspace --- the simplex laws of
Section~\ref{sec:cmh-exact-cases} are the case of interest --- $\Sigma^{-1}$ denotes the inverse
on that affine tangent space, equivalently the Moore--Penrose inverse in ambient coordinates.
\end{definition}

Writing $K=-\Div_\mu(H\Sigma^{-1}H\nabla\,\cdot\,)$, the definition is the quadratic-form
comparison $K\preceq C\Aop^2$, equivalently the second-order Riesz-transform bound
\begin{equation}
  \norm{\Sigma^{-1/2}H\nabla\Aop^{-1}}_{L^2(\mu)\to L^2(\mu;\R^n)}\le\sqrt C.
  \label{eq:cmh-riesz}
\end{equation}
This is the precise content of the schema \eqref{eq:cmh4-schema}: $\Sigma$ is the covariance,
$L$ is the Stein generator \eqref{eq:stein-generator}, the measure is $\mu$ itself, and the
admissible class is $\Dom(\Aop)$.

\subsection{CMH implies the affine Poincar\'e inequality}
\label{subsec:cmh-implies}

\begin{theorem}[Endpoint reduction; answers Question~\ref{q:cmh-normalization}]
\label{thm:cmh-implies-affine-poincare}
For every centered log-concave $\mu$ in the regular moment-map class described above,
\begin{equation}
  \CPaff(\mu)\le\CMH(\mu).
  \label{eq:cmh-implies-poincare}
\end{equation}
In particular a universal bound $\CMH\le4$ over all isotropic log-concave measures in all
dimensions would prove Conjecture~\ref{conj:kls} with affine Poincar\'e constant $4$.
\end{theorem}

\begin{proof}
Put $C=\CMH(\mu)$; there is nothing to prove if $C=\infty$.  Let
$\mathscr C$ be the restrictions to $\operatorname{supp}\mu$ of
$\R+C_c^\infty(\R^n)$ and first take a centered $f\in\mathscr C$.  This class lies in both
form domains even when $H$ is unbounded: the Stein normalization $\E_\mu H=\Sigma$ gives
\[
  \E_\mu\inner{H\nabla f}{\nabla f}
  \le \norm{\nabla f}_\infty^2\E_\mu\Tr H
  =\norm{\nabla f}_\infty^2\Tr\Sigma<\infty.
\]
For
$0<\eps<R$ set
\[
  \Pi_{\eps,R}=\one_{[\eps,R]}(\Aop),
  \qquad g_{\eps,R}=\Pi_{\eps,R}\Aop^{-1}f\in\Dom(\Aop).
\]
Then $\Aop g_{\eps,R}=\Pi_{\eps,R}f$.  The form--operator pairing, rather than any commutation
of $\Pi_{\eps,R}$ with the $\Sigma$-form, gives
\begin{align*}
  \norm{\Pi_{\eps,R}f}_2^2
  &=\inner{f}{\Pi_{\eps,R}f}_{L^2}
    =\inner{f}{\Aop g_{\eps,R}}_{L^2}
    =\E_\mu\inner{\nabla f}{H\nabla g_{\eps,R}}\\
  &\le
    \bigl(\E_\mu\inner{\Sigma\nabla f}{\nabla f}\bigr)^{1/2}
    \bigl(\E_\mu\inner{H\nabla g_{\eps,R}}
                            {\Sigma^{-1}H\nabla g_{\eps,R}}\bigr)^{1/2}\\
  &\le C^{1/2}
    \bigl(\E_\mu\inner{\Sigma\nabla f}{\nabla f}\bigr)^{1/2}
    \norm{\Pi_{\eps,R}f}_2.
\end{align*}
After division (with the zero case harmless) and squaring,
\[
  \norm{\Pi_{\eps,R}f}_2^2
  \le C\E_\mu\inner{\Sigma\nabla f}{\nabla f}.
\]
Because $H\succ0$ and the support is connected, $\ker\Aop$ consists of the constants.  Thus
$f\perp\ker\Aop$, and the spectral theorem gives
$\Pi_{\eps,R}f\to f$ in $L^2$ as $R\to\infty$ and then $\eps\downarrow0$.

Finally define $H^1_\Sigma(\mu)$ as the closure of $\mathscr C$ for
$\norm{f}_2^2+\E_\mu\inner{\Sigma\nabla f}{\nabla f}$ (equivalently the usual $H^1(\mu)$ here,
since $\Sigma\succ0$ is constant).  Approximate an arbitrary $f\in H^1_\Sigma(\mu)$ by
functions in $\mathscr C$ and subtract their means.  Both the variances and the
$\Sigma$-energies converge, so
the preceding inequality passes to the limit.  Notice that this density step never asserts
that finite $\Sigma$-energy implies finite $H$-energy when $H$ is unbounded.
\end{proof}

The proof uses no Bochner identity, no positivity of the Monge--Amp\`ere remainder, and no
property of $H$ beyond symmetry, positivity, and $\Div_\mu H=-p$. That economy is the reason
the constant transfers with no loss: the reduction is exactly one Cauchy--Schwarz.

\subsection{The Hodge content: the affine channel and the solenoidal excess}
\label{subsec:cmh-hodge}

Set $u=H\nabla g$ and $h=-\Div_\mu u=-L_\mu g$. Let $\psi$ solve the ordinary weighted
covariance divergence problem $-\Div_\mu(\Sigma\nabla\psi)=h$ on the centered subspace, and put
$w=u-\Sigma\nabla\psi$, so that $\Div_\mu w=0$.

\begin{proposition}[Weighted Hodge identity]
\label{prop:cmh-hodge}
The fields $\Sigma\nabla\psi$ and $w$ are orthogonal in $L^2(\mu;\Sigma^{-1})$, and
\begin{equation}
  \E_\mu\inner{u}{\Sigma^{-1}u}
  =\E_\mu\inner{\Sigma\nabla\psi}{\nabla\psi}+\E_\mu\inner{w}{\Sigma^{-1}w}.
  \label{eq:cmh-hodge}
\end{equation}
Moreover the operator norm of $h\mapsto\E_\mu\inner{\Sigma\nabla\psi}{\nabla\psi}$ on the
centered subspace of $L^2(\mu)$ is exactly $\CPaff(\mu)$.
\end{proposition}

\begin{proof}
Weighted integration by parts gives
$\E\inner{\Sigma\nabla\psi}{\Sigma^{-1}w}=\E\inner{\nabla\psi}{w}=-\E[\psi\Div_\mu w]=0$, and
expanding the square proves \eqref{eq:cmh-hodge}. For the last claim, let
$\Aop_1=-\Div_\mu(\Sigma\nabla\,\cdot\,)$, so $\psi=\Aop_1^{-1}h$ and
$\E\inner{\Sigma\nabla\psi}{\nabla\psi}=\inner{h}{\Aop_1^{-1}h}$; the supremum over $h$ of the
ratio to $\norm{h}_2^2$ is $\norm{\Aop_1^{-1}}$, which is $\CPaff(\mu)$ by
\eqref{eq:affine-poincare-constant}.
\end{proof}

\begin{corollary}[The exact Hodge consequence]
\label{rem:cmh-stronger-than-kls}
$\Sigma\nabla\psi$ is the least-$L^2(\Sigma^{-1})$ field with divergence $-h$, and by
Proposition~\ref{prop:cmh-hodge} controlling it \emph{is} the affine Poincar\'e inequality. CMH
demands in addition that the solenoidal excess $w$ be paid for within the same budget.

In dimension one a square-integrable divergence-free field with the no-flux convention of
\S\ref{subsec:cmh-conventions} vanishes identically, so $w=0$ and CMH degenerates exactly to
the Poincar\'e inverse-divergence problem --- this is why Theorem~\ref{thm:cmh-1d} is an
identity rather than an inequality.  In every dimension the decomposition proves the exact
comparison
\[
  \CMH(\mu)\ge\CPaff(\mu).
\]
In dimension at least two the divergence-free subspace is nontrivial, so the identity displays
an additional channel that CMH must control.
\end{corollary}

\begin{remark}[What the Hodge identity does not prove]
The identity does \emph{not} show that this channel
is nonzero for a CMH extremizing sequence, nor does it exhibit a measure separating CMH from
$\CPaff$.  Thus $\mathrm{CMH}(4)$ is sufficient for KLS with constant $4$, while equivalence
and strict nonimplication are both open.  The ledger records $\mathrm{CMH}(4)$ as a route target,
not as an established reformulation or a proved strict strengthening of KLS.
\end{remark}

The product formula below shows that products of one-sided exponentials saturate
$\mathrm{CMH}(4)$ with zero slack.  Corollary~\ref{rem:cmh-saturation-risk} therefore turns the
possible solenoidal gap into a concrete perturbative test, but not into a proved separation.

\subsection{Gate zero: the necessary linear-sector condition}
\label{subsec:gate-zero}

Take $g(p)=a\cdot p$ in \eqref{eq:cmh-constant}. Then $D^2g=0$, $L_\mu g=-a\cdot p$, so
$\E(L_\mu g)^2=a^\top\Sigma a$ while the numerator is $a^\top\E[H\Sigma^{-1}H]a$. Hence:

\begin{conjecture}[Gate zero]
\label{conj:gate-zero}
Every centered log-concave moment measure satisfies
\begin{equation}
  \E\bigl[H\Sigma^{-1}H\bigr]\preceq4\Sigma,
  \label{eq:gate-zero}
\end{equation}
equivalently $\E H^2\preceq4\Id$ in isotropic position.
\end{conjecture}

This is necessary for universal $\mathrm{CMH}(4)$ and is not a known consequence of KLS. It is
the cheapest falsifiable consequence of the whole route: a single matrix expectation, with no
test function and no operator inverse. A measure with
$\lmax(\Sigma^{-1/2}\E[H\Sigma^{-1}H]\Sigma^{-1/2})>4$ would disprove $\mathrm{CMH}(4)$ without
disproving Conjecture~\ref{conj:kls}.

\begin{corollary}[Gate zero is a trace-upgrade problem]
\label{rem:gate-zero-trace-upgrade}
Theorem~\ref{thm:chen-klartag-moment-hessian} gives
$\E_\nu\norm{H}_{\HS}^2\le2n$, that is $\Tr(\E H^2)\le2n$: the \emph{average} eigenvalue of
$\E H^2$ is already at most $2$ in isotropic position. Gate zero asks for the \emph{largest}
eigenvalue to be at most $4$. It is therefore an operator-to-trace upgrade carrying a factor-$2$
budget.
\end{corollary}

\begin{remark}[Relation to the shared trace-upgrade difficulty]
Gate zero belongs to the same difficulty family as \texttt{q:upgrade}, the high-rank part of
\texttt{q:stein-weighted}, and \texttt{q:alignment}; no equivalence is asserted, exactly as in
Remark~\ref{rem:trace-upgrade-unification}. The practical consequence is a split verdict:
gate zero is cheap to \emph{test} on a model and is expected to be as hard to \emph{prove} as
the rest of the program. Only the first half is dispatchable, and
\texttt{experiments/README.md} records the corresponding numerical channel.
\end{remark}

\subsection{Why the constant-matrix estimate cannot supply gate zero}
\label{subsec:gate-zero-countermodel}

In isotropic position define the positive self-adjoint superoperator $\calT(B)=\E[HBH]$ on
symmetric matrices with the Hilbert--Schmidt inner product.
Theorem~\ref{thm:letwin-moment-map} says $\calT\preceq2\,\Id$; gate zero asks instead for
$\calT(\Id)=\E H^2\preceq4\,\Id$. The two differ by the order of the noncommuting factors, and
the gap is exactly one commutator:
\begin{equation}
  \Tr(B^2H^2)=\Tr(BHBH)+\tfrac12\norm{[B,H]}_{\HS}^2
  \qquad(B,H\ \text{symmetric}).
  \label{eq:static-commutator}
\end{equation}
Indeed $[B,H]$ is antisymmetric, so
$\norm{[B,H]}_{\HS}^2=-\Tr([B,H]^2)=2\Tr(B^2H^2)-2\Tr(BHBH)$. The first term of
\eqref{eq:static-commutator} is what Letwin's theorem controls; the second is the \emph{static
transverse commutator}, and the following shows it is genuinely unconstrained by matrix moments.

\begin{proposition}[Algebraic countermodel]
\label{prop:letwin-not-gate-zero}
For every $m\ge18$ there is a random positive semidefinite $(m+1)\times(m+1)$ matrix $H$ with
$\E H=\Id$ and
\[
  \E\Tr(BHBH)\le2\Tr(B^2)\quad\text{for every symmetric }B,
\]
yet $\lmax(\E H^2)>4$.
\end{proposition}

\begin{proof}
Let $z$ be uniform on $S^{m-1}$ and set
\[
  d=\frac m{\sqrt{2m-1}},\qquad c=1-\frac dm,\qquad
  H(z)=\begin{pmatrix}1&\sqrt d\,z^\top\\[2pt]\sqrt d\,z&c\,\Id_m+d\,zz^\top\end{pmatrix}.
\]
The Schur complement of the upper-left entry is $c\,\Id_m\succ0$ because $c>0$, so $H\succeq0$;
since $\E z=0$ and $\E zz^\top=\Id_m/m$ we get
$\E H=1\oplus(c+d/m)\,\Id_m=\Id_{m+1}$.

Write $B=\begin{psmallmatrix}a&r^\top\\r&D\end{psmallmatrix}$ with $D$ symmetric. Using
$\E(z^\top Dz)^2=\bigl((\Tr D)^2+2\Tr(D^2)\bigr)/(m(m+2))$,
\begin{align*}
  \E\Tr(BHBH)
  &=a^2+\frac{2d}ma\Tr D+2\Bigl(c+\frac{2d}m\Bigr)\abs r^2\\
  &\quad+\Bigl(c^2+\frac{2cd}m+\frac{2d^2}{m(m+2)}\Bigr)\Tr(D^2)
   +\frac{d^2}{m(m+2)}(\Tr D)^2.
\end{align*}
Writing $D=t\,\Id_m+D_0$ with $\Tr D_0=0$ splits the form into three $O(m)$-invariant sectors
that do not interact: the cross-vector sector, the traceless-$D$ sector, and the two-dimensional
scalar sector $(a,t)$. Comparing with
$2\Tr(B^2)=2a^2+4\abs r^2+2\Tr(D^2)$ sector by sector:
the vector coefficient is $2(c+2d/m)=2(1+d/m)\le4$ since $d\le m$; the traceless coefficient is
$c^2+2cd/m+2d^2/(m(m+2))=1+(m-2)/((2m-1)(m+2))\le2$; and on the scalar sector, using
$mc^2+2cd+d^2=m+d^2-d^2/m$, the deficit is
\[
  2(a^2+mt^2)-\E\Tr(BHBH)=a^2-2dat+\Bigl[m-d^2\Bigl(1-\frac1m\Bigr)\Bigr]t^2=(a-dt)^2\ge0,
\]
the last equality because $d^2(2-1/m)=m$ by the choice $d^2=m^2/(2m-1)$. Hence
$\E\Tr(BHBH)\le2\Tr(B^2)$ for every symmetric $B$, with equality on the scalar ray $a=dt$.

Finally $e_1^\top\E H^2e_1=1+d\abs z^2=1+d$, and $1+d>4$ is equivalent to $m^2>9(2m-1)$, i.e.
$m^2-18m+9>0$, which holds precisely for $m\ge18$.
\end{proof}

\begin{remark}[Scope of the countermodel]
\label{rem:countermodel-scope}
The matrices above are \emph{not} claimed to be moment-map Hessians; no Monge--Amp\`ere,
Codazzi, or Hessian-compatibility condition is imposed. The proposition proves only that
positivity, the normalization $\E H=\Id$, and the constant-matrix estimate
\eqref{eq:letwin-matrix} do not imply gate zero by matrix algebra. Any proof of
\eqref{eq:gate-zero} must therefore consume differential moment-map structure. Equivalently, by
\eqref{eq:static-commutator}, it must control $\E\norm{[B,H]}_{\HS}^2$, and
Proposition~\ref{prop:letwin-not-gate-zero} exhibits an admissible law where that quantity is
maximal on the scalar ray. This is the same obstruction as
Question~\ref{q:mm-square-root-commutator} in its most elementary, static,
finite-dimensional form.
\end{remark}

\begin{question}[Gate zero on genuine moment maps]
\label{q:gate-zero}
Decide \eqref{eq:gate-zero}. Either prove
$\E[H\Sigma^{-1}H]\preceq4\Sigma$ for every log-concave moment measure using differentiated
Monge--Amp\`ere structure --- which by Remark~\ref{rem:countermodel-scope} means a
dimension-free bound on the static commutator $\E\norm{[B,H]}_{\HS}^2$ for the relevant
multiplier class --- or exhibit a genuine moment map with
$\lmax(\Sigma^{-1/2}\E[H\Sigma^{-1}H]\Sigma^{-1/2})>4$, which refutes $\mathrm{CMH}(4)$ without
refuting Conjecture~\ref{conj:kls}.
\end{question}

\subsection{Conventions, domains, and affine covariance}
\label{subsec:cmh-conventions}

\paragraph{Weighted divergence and boundary flux.}
$\Div_\mu u=\rho^{-1}\Div(\rho u)$. All integrations by parts are justified first for compactly
supported fields or fields with vanishing normal flux, and the closed forms are obtained by
completion. The no-flux convention is essential in dimension one, where it is what rules out a
nonzero constant divergence-free field and forces $w=0$ in
Corollary~\ref{rem:cmh-stronger-than-kls}.

\paragraph{Closed operators.}
$\calE_H(f,g)=\E_\mu\inner{H\nabla f}{\nabla g}$ is closable under the smooth moment-map
hypotheses of \S\ref{subsec:mm-coordinates}; $\Aop=-L_\mu$ is its nonnegative self-adjoint
operator, and $\Aop^{-1}$ always means the pseudoinverse on $(\ker\Aop)^\perp$. The spectral
truncation in the proof of Theorem~\ref{thm:cmh-implies-affine-poincare} avoids assuming a
spectral gap.

\paragraph{Affine covariance.}
If $T$ is invertible and $Y=TX$, then $\Sigma_Y=T\Sigma_XT^\top$ and
$H_Y(Tx)=TH_X(x)T^\top$; with $g_Y(y)=g_X(T^{-1}y)$ both the numerator and denominator of
\eqref{eq:cmh-constant} are unchanged. Hence $\CMH$ is an invariant of the affine equivalence
class, matching the affine invariance of $\CPaff$. For a \emph{noninvertible} map the
transported Stein kernel need not be the canonical moment-map kernel of the image, so only the
Poincar\'e-level statement of Corollary~\ref{cor:cmh-linear-images} is available there.

\paragraph{Conditional closure for general log-concave measures.}
Question~\ref{q:cmh-approximation} is certified at the Poincar\'e level as follows. For an
arbitrary centered log-concave $\mu$, centered Gaussian-convolution, Gaussian-tilt, and
growing-ball approximants $\mu_q$ belong to the published compact-target regular class of
Theorem~\ref{thm:regular-moment-map-compact-target} and converge to $\mu$ in ambient $W_2$ with
second moments, hence $\Sigma_q\to\Sigma$. If $\sup_q\CMH(\mu_q)\le C$, then for every smooth
compactly supported $f$,
\[
  \Var_{\mu_q}f\le C\int\inner{\Sigma_q\nabla f}{\nabla f}\,d\mu_q.
\]
Both sides converge because $f$ and $\nabla f$ are bounded and continuous and the covariances
converge. Intrinsic smooth-core density on $S=\operatorname{Ran}\Sigma$ extends the inequality to
the closed covariance-form relaxation $H^1_\Sigma(\mu)$, while $\Sigma$ and $\Sigma^+$ annihilate
ambient normal derivatives. This proves constant-preserving closure, including singular
covariance. It does not prove the uniform premise and asserts no continuity of $\CMH$.
\end{document}
